\chapter{Limits and Continuity of Functions}
\chaptermark{Limits and Continuity}
\label{ch:continuity}

\section{Limits at a Point}

Sequence limits describe processes indexed by natural numbers.  A function
limit asks the same question when the input can approach a point through all
nearby domain values, so the value assigned at the target itself must be
kept logically separate.

\begin{definition}[Limit point]
Let $D\subseteq\R$ and $a\in\R$. The point $a$ is a \emph{limit point} of $D$ if, for every $\delta>0$, there is an $x\in D$ such that $0<\abs{x-a}<\delta$.
\end{definition}

There are points of $D$ other than $a$ arbitrarily close to $a$.
Membership and approachability are different questions. For $D=(0,1)$,
the point $0$ is a limit point although $0\notin D$: for any $\delta>0$,

\[
x=\min\{\delta/2,1/2\}
\]
 belongs to $D$ and satisfies $0<x<\delta$.
Thus a right-hand limit at this missing endpoint is meaningful.
For $D=\{0\}\cup[1,2]$, however, $0\in D$ is not a limit point:
no point of $D$ satisfies
\[
0<|x|<1/2.
\]


\begin{definition}[Isolated point]
A point $a\in D$ is \textbf{isolated in $D$} if some $\delta>0$ satisfies
$D\cap(a-\delta,a+\delta)=\{a\}$. For $a\in D$, negating the
limit-point definition shows that being isolated is equivalent to not
being a limit point of $D$.
\end{definition}

\begin{proposition}[Sequential test for a limit point]
\label{prop:limit-point-sequence}
The point $a$ is a limit point of $D$ if and only if there is a sequence
$x_n\in D\setminus\{a\}$ with $x_n\to a$.
\end{proposition}
\begin{proof}
Choose
\[
0<|x_n-a|<1/n
\]
 in one direction. Conversely, convergence puts
some $x_n\ne a$ in every prescribed punctured neighborhood of $a$.
\end{proof}
We can therefore test function limits by sequences; the
limit-point hypothesis guarantees that there really are approaching inputs.

\begin{definition}[Function limit]
\label{def:function-limit}
Let $f:D\to\R$, let $a$ be a limit point of $D$, and let $L\in\R$. We write
\[
\lim_{x\to a}f(x)=L
\]
if, for every $\eps>0$, there is $\delta>0$ such that every $x\in D$ satisfying $0<\abs{x-a}<\delta$ also satisfies $\abs{f(x)-L}<\eps$.
\end{definition}

\paragraph{Why require a limit point?}
If $a$ is not a limit point, some punctured neighborhood contains no domain
points. The implication in the limit definition would then be true for
\emph{every} proposed $L$, because there is no admissible $x$ to test.
It could not determine a unique value. In the uniqueness proof below,
the limit-point assumption is exactly what supplies an actual $x\ne a$
close enough to compare both proposed limits. Notice that $f(a)$ is not
needed and, when defined, has no effect on the punctured limit.

\paragraph{Discover backward; prove forward.}
The target $|f(x)-L|<\eps$ suggests a sufficient restriction on
$|x-a|$. During discovery we may start at the target and work backward.
The proof must start with a positive $\delta$ chosen from $\eps$ and
then verify the implication for \emph{every} admissible $x$.
For $3x-1\to5$ at $2$, discovery gives $3|x-2|<\eps$;
the forward proof chooses
\[
\delta=\eps/3
\]
 and checks
$0<|x-2|<\delta\Rightarrow |3x-6|<3\delta=\eps$.
The next example adds a local bound before allocating the error.

\begin{example}[Localize before choosing the tolerance]
To prove $x^2\to a^2$, discovery begins with
$|x^2-a^2|=|x-a||x+a|$. The second factor is not fixed. First impose
$|x-a|<1$, which gives $|x+a|<1+2|a|$. Now choose
\[
\delta=\min\left\{1,\frac{\eps}{1+2|a|}\right\}.
\]
For $0<|x-a|<\delta$, the product is less than $\eps$. The pattern is to bound
the uncontrolled factor locally and only then choose the accuracy.
\end{example}

\begin{example}[Controlling a denominator directly]
To prove
\[
1/x\to1/2
\]
 as $x\to2$, first require $|x-2|<1$.
Then $x>1$ and

\[
|1/x-1/2|=|x-2|/(2|x|)\leq|x-2|/2.
\]

Given $\eps>0$, choose $\delta=\min\{1,2\eps\}$. For every
$0<|x-2|<\delta$ the preceding estimate is strictly less than $\eps$.
A restriction that keeps a denominator away from zero is part of the
proof, not an optional algebraic convenience.
\end{example}

\begin{theorem}[Uniqueness of function limits]
\label{thm:unique-function-limit}
Under the hypotheses of \cref{def:function-limit}, there is at most one value $L$ for which the displayed limit holds.
\end{theorem}
\begin{proof}
Suppose $L$ and $M$ both satisfy the definition. Given $\eps>0$, obtain positive numbers $\delta_L,\delta_M$ for tolerance
\[
\eps/2.
\]
 Since $a$ is a limit point, choose $x\in D$ with
\[
0<\abs{x-a}<\min\{\delta_L,\delta_M\}.
\]
Then
\[
\abs{L-M}\leq\abs{L-f(x)}+\abs{f(x)-M}<\eps.
\]
By \cref{lem:positive-small}, $L=M$.
\end{proof}

\begin{theorem}[Sequential criterion]
\label{thm:sequential-function-limit}
Let $f:D\to\R$ and let $a$ be a limit point of $D$. Then
\[
\lim_{x\to a}f(x)=L
\]
if and only if, for every sequence $(x_n)$ in $D\setminus\{a\}$ with $x_n\to a$, one has $f(x_n)\to L$.
\end{theorem}
\begin{proof}
Suppose first that the function limit exists. Given $\eps>0$, take $\delta$ from \cref{def:function-limit}. Choose $N$ with $|x_n-a|<\delta$ for all $n\geq N$.
For every such $n$, $x_n\ne a$, so $|f(x_n)-L|<\eps$.

Conversely, suppose the sequential statement holds but the function limit
fails. Negating the definition gives
\[
\exists\eps_0>0\ \forall\delta>0\ \exists x\in D:
\quad 0<|x-a|<\delta\quad\text{and}\quad |f(x)-L|\geq\eps_0.
\]
The order of quantifiers gives one fixed $\eps_0$ for all choices of
$\delta$. For each $n$, set
\[
\delta=1/n
\]
 and choose the corresponding
point $x_n\in D$, so
\[
0<\abs{x_n-a}<\frac1n
\qquad\text{and}\qquad
\abs{f(x_n)-L}\geq\eps_0.
\]
For any $\eta>0$, choose
\[
N>1/\eta;
\]
 then $n\geq N$ implies

\[
|x_n-a|<1/n\leq1/N<\eta.
\]
 Hence $x_n\to a$, while the second prevents $f(x_n)$ from converging to $L$, a contradiction.
\end{proof}

\begin{definition}[One-sided and infinite limits]
Let $f:D\to\R$. For a right-hand limit assume that $a$ is a limit point
 of $D\cap(a,\infty)$; use $D\cap(-\infty,a)$ for a left-hand limit. The assertion
\[
\lim_{x\to a^+}f(x)=L
\]
means that for every $\eps>0$ there is $\delta>0$ such that
\[
x\in D,\quad 0<x-a<\delta
\quad\Longrightarrow\quad
\abs{f(x)-L}<\eps.
\]
The left-hand limit is defined with $0<a-x<\delta$.  We write
$f(x)\to+\infty$ as $x\to a^+$ if for every $M\in\R$ there is
$\delta>0$ such that the same hypotheses imply $f(x)>M$; replace the last
inequality by $f(x)<M$ for a limit of $-\infty$.  The analogous two-sided
definitions use $0<\abs{x-a}<\delta$.

At an unbounded end of the domain (assume $D$ is unbounded in the stated direction),
\[
\lim_{x\to+\infty}f(x)=L
\]
means that for every $\eps>0$ there is $A\in\R$ such that $x\in D$ and
$x>A$ imply $\abs{f(x)-L}<\eps$.  Replacing this conclusion by $f(x)>M$
defines $f(x)\to+\infty$ as $x\to+\infty$; changing the direction of the
input or output inequalities gives the other cases.
\end{definition}

\begin{remark}[Infinity is behavior, not a value]
The notation $f(x)\to+\infty$ does not assert convergence to a member of
$\R$.  It abbreviates an eventual-inequality statement.  For example,

\[
1/x\to+\infty
\]
 as $x\to0^+$, but the two-sided limit at $0$ does not tend
to either sign of infinity because
\[
1/x\to-\infty
\]
 from the left.
\end{remark}

\begin{theorem}[Algebraic function-limit laws]
Let $a$ be a limit point of a common domain $D$, and suppose
$f(x)\to L$ and $g(x)\to M$ as $x\to a$ through $D$.
For constants $c,d$, $cf+dg\to cL+dM$, $fg\to LM$, and if
$M\ne0$,
\[
f/g\to L/M
\]
 on a sufficiently small punctured domain.
\end{theorem}
\begin{proof}
For every $x_n\in D\setminus\{a\}$ tending to $a$, transfer the sequence
laws through the sequential criterion. The hypotheses give
$f(x_n)\to L$, $g(x_n)\to M$. The sequence sum
and product laws give the asserted limits on every such sequence, hence
the function limits. If $M\ne0$, choose $\delta_0>0$ such that
$0<|x-a|<\delta_0$ implies
\[
|g(x)-M|<|M|/2;
\]
 then
\[
|g(x)|>|M|/2
\]
 and the quotient is defined locally. Apply the
sequence quotient law and the same criterion on that punctured domain.
\end{proof}
\begin{example}[Direct proof and rationalization]
To prove
\[
\lim_{x\to2}(3x-1)=5,
\]
observe that $\abs{(3x-1)-5}=3\abs{x-2}$ and choose

\[
\delta=\eps/3.
\]
  For a less direct example,
\[
\lim_{x\to0}\frac{\sqrt{1+x}-1}{x}
=\lim_{x\to0}\frac1{\sqrt{1+x}+1}=\frac12,
\]
where the equality for $x\ne0$ follows by rationalization and the last
step follows directly from an estimate, without invoking continuity.
For
\[
|x|<1/2,
\]
 the positive root is defined by the real-number root
construction, and
\[
\left|\frac1{\sqrt{1+x}+1}-\frac12\right|
=\frac{|x|}{2(\sqrt{1+x}+1)^2}\leq\frac{|x|}{2}.
\]
Choosing
\[
\delta=\min\{1/2,2\eps\}
\]
 proves the limit.
\end{example}

\begin{example}[One-sided limits detect nonexistence]
For
\[
f(x)=x/\abs{x}
\]
 on $\R\setminus\{0\}$,
\[
\lim_{x\to0^-}f(x)=-1,
\qquad
\lim_{x\to0^+}f(x)=1.
\]
Hence the two-sided limit does not exist.  Indeed, the sequences

\[
x_n=-1/n,\qquad y_n=1/n
\]
 approach the same point while their images have
different limits.
\end{example}

For one-sided approach, $t\downarrow a$ means $t\to a$ from above
($t\to a^+$), while $t\uparrow a$ means $t\to a$ from below
($t\to a^-$). This related notation concerns a variable approaching a point;
it is grammatically distinct from the sequence notation $a_n\uparrow L$.

\begin{remark}[A working strategy for function limits]
First discover a plausible answer by substitution, simplification, a graph,
or comparison.  Next prove it: use limit laws only after their hypotheses
are met; cancel only on the punctured neighborhood; rationalize radicals;
or bound the error by a simpler expression and squeeze.  For nonexistence,
compare one-sided limits or produce two approaching sequences with
different image limits.  The exact value $f(a)$ is irrelevant to
$\lim_{x\to a}f(x)$.
\end{remark}

\begin{example}[A zero denominator is a question, not an answer]
At $x=1$, the quotient
\[
(x^2-1)/(x-1)
\]
 is undefined, but on the
punctured domain it equals $x+1$ and its limit is $2$.
In contrast
\[
1/x
\]
 has incompatible infinite limits at zero, while

\[
x^2/x=x
\]
 has limit zero. In all three cases the denominator tends
to zero, so the ordinary quotient theorem cannot decide the problem.
For a direct disproof that
\[
x/|x|\to1
\]
 at zero, take $\eps_0=1$.
For every $\delta>0$,
\[
x=-\delta/2
\]
 has $0<|x|<\delta$ and

\[
|x/|x|-1|=2\geq\eps_0.
\]
 To exclude all proposed limits, use the
two approaching sequences in the preceding example.
\end{example}

\subsection*{Polynomial and rational limits}
The identity function tends to $a$ as $x\to a$, directly with
$\delta=\eps$. Repeated products and sums in the limit laws therefore give
\[
P(x)=a_nx^n+\cdots+a_0
\quad\Longrightarrow\quad\lim_{x\to a}P(x)=P(a).
\]
If $Q(a)\ne0$, denominator control and the quotient law give

\[
P(x)/Q(x)\to P(a)/Q(a).
\]
 For example,

\[
(x^2+1)/(x+2)\to5/4
\]
 as $x\to2$.

The same algebraic laws hold at infinity. To justify this directly,
in the sequence-limit proofs replace every tail condition $n\geq N_i$
by $x>A_i$ and use $A=\max_i A_i$; the estimates do not otherwise
change. The basic input is
\[
1/x\to0:
\]
 choose

\[
A=\max\{1,1/\eps\},
\]
 so $x>A$ gives
\[
|1/x|<\eps.
\]

Integer powers and limit arithmetic handle all lower-order terms.

For the expressions $f(x)=a_nx^n+\cdots+a_0$ and
$g(x)=b_mx^m+\cdots+b_0$, with nonzero leading coefficients and
degrees $n,m$, divide by the dominant powers:
\[
F(x)=\frac{f(x)}{x^n}\to a_n,
\qquad
G(x)=\frac{g(x)}{x^m}\to b_m.
\]
Choose $A_0$ with
\[
|G(x)-b_m|<|b_m|/2
\]
 for $x>A_0$.
Then
\[
|G(x)|>|b_m|/2,,\qquad H(x)=F(x)/G(x)\to c=a_n/b_m.
\]

Consequently
\[
f(x)/g(x)=x^{n-m}H(x)
\]
 tends to $0$ if $n<m$,
to $c$ if $n=m$, and to the infinity with sign $c$ if $n>m$.
For completeness, when $c>0$ choose $A_1$ with
\[
H(x)>c/2
\]
 for
$x>A_1$. For any target $M$, every
\[
x>A=\max\left\{A_0,A_1,1,\frac{2\max\{M,0\}}c\right\}
\]
satisfies
\[
x^{n-m}H(x)>(c/2)x>M.
\]
 Negating handles $c<0$.
This is the mechanism of \cref{prop:polynomial-sequence}, now for
genuine function limits, not just natural-number substitutions.

For $x\to-\infty$, put $t=-x\to+\infty$.
The leading coefficients become $(-1)^na_n$ and $(-1)^mb_m$.
The limit is still zero for $n<m$ and $c$ for $n=m$;
for $n>m$ its infinite sign is that of $(-1)^{n-m}c$.
For example
\[
(2x^3+1)/(x+1)
\]
 tends to $+\infty$ at both ends,
whereas
\[
(2x^2+1)/(x+1)
\]
 tends to opposite infinities.

\section{Asymptotic Notation}

Limit arguments often need only a comparison of sizes rather than an exact
formula. We use the following notation for two common kinds of
comparison without replacing the underlying estimates.

\begin{definition}[Big-O and little-o]
Let $f,g$ be real-valued functions defined sufficiently near $a$.  We write
\[
f(x)=O(g(x))\qquad(x\to a)
\]
if there are $C>0$ and a punctured neighborhood of $a$ on which
\[
\abs{f(x)}\leq C\abs{g(x)}.
\]
We write
\[
f(x)=o(g(x))\qquad(x\to a)
\]
if $g(x)\ne0$ sufficiently near $a$ and
\[
\frac{f(x)}{g(x)}\longrightarrow0.
\]
For sequences, $a_n=O(b_n)$ as $n\to\infty$ and $a_n=o(b_n)$ as
$n\to\infty$ have the identical eventual-bound and quotient-limit meanings.
\end{definition}

The limiting regime is part of the assertion. For example,
$x^2=o(x)$ as $x\to0$, but not as $x\to\infty$; and
$n^{-2}=o(n^{-1})$ as $n\to\infty$. All rules below refer to one fixed
regime, with comparison functions eventually nonzero when quotients occur.

\paragraph{Translating the notation.}
The assertion $r=o(g)$ means $r=\epsilon(x)g(x)$ with
$\epsilon(x)\to0$. Equivalently, for every $\eps>0$,
$|r(x)|\leq\eps|g(x)|$ sufficiently near the limit. The assertion
$r=O(g)$ uses one fixed constant $C$, rather than every positive tolerance.
When $g$ is eventually nonzero it is equivalent to $r=C(x)g(x)$ with
$C(x)$ bounded. These are functions obeying estimates, not mysterious
individual objects called $O(g)$ or $o(g)$. Thus $f=g+o(g)$ means that
there exists a remainder $r=o(g)$ for which $f=g+r$. Different occurrences
of $o(g)$ may denote different remainders.

\begin{proposition}[When does a remainder vanish?]
\label{prop:asymptotic-vanishing}
One has $r=o(1)$ exactly when $r\to0$. The assertion $r=O(1)$ says
only that $r$ is bounded. If $g\to0$, either $r=O(g)$ or $r=o(g)$
implies $r\to0$.
\end{proposition}
\begin{proof}
The first two statements translate the definitions with comparison scale
$1$. For big-O use $|r|\leq C|g|$ and squeeze. Little-o implies big-O
by choosing tolerance $1$, so the same argument applies.
\end{proof}
For example $\sin x=O(1)$ as $x\to\infty$ but has no limit (use

\[
x=2\pi n+\pi/2,\qquad x=2\pi n+3\pi/2
\]
). As $x\to\infty$,
$x=o(x^2)$ and $\sqrt{x}=o(x)$, although both remainders diverge.
Little-o means small \emph{relative to a scale}, not necessarily small
in absolute size.
For vector-valued remainders the same definitions use the norm of the
remainder; Chapter~10 proves the norm estimates needed for composition.

Near $h=0$, the three conditions are strictly ordered:
\[
r(h)=o(|h|)\ \Longrightarrow\ r(h)=O(|h|)\ \Longrightarrow\ r(h)\to0.
\]
Neither converse holds: $|h|$ is $O(|h|)$ but not $o(|h|)$, while
$\sqrt{|h|}\to0$ but is not $O(|h|)$. In contrast, $|h|^2=o(|h|)$.
Chapter~10 replaces absolute values by norms with the same examples.

\begin{proposition}[A calculus of remainder estimates]
\label{prop:asymptotic-calculus}
In a fixed regime,
\[
\begin{gathered}
O(g)+O(g)=O(g),\qquad o(g)+o(g)=o(g),\\
O(f)O(g)=O(fg),\qquad o(f)O(g)=o(fg),\qquad o(f)o(g)=o(fg).
\end{gathered}
\]
Also $r=O(f), f=O(g)$ imply $r=O(g)$; replacing either of these
two hypotheses by little-o gives $r=o(g)$.
\end{proposition}
\begin{proof}
Translate each remainder into its comparison scale times a
bounded or vanishing coefficient. For sums the coefficients add; for
products they multiply. Sums of bounded coefficients are bounded, sums
of vanishing coefficients vanish, and a vanishing coefficient times a
bounded one vanishes. A vanishing coefficient is itself eventually bounded.
For transitivity, multiply
\[
r/f
\]
 by
\[
f/g;
\]
 these same facts apply.
For big-O sums even at zeros of $g$, the triangle inequality gives
$|r_1+r_2|\leq(C_1+C_2)|g|$; product bounds work in the same way.
\end{proof}
These displayed ``equalities'' assert estimates for arbitrary remainders;
they do not license cancellation as if $O(g)$ were a fixed number.
In particular, $O(g)-O(g)$ need not be zero, and the reciprocal of an
$O(1)$ function need not be $O(1)$ without a bound away from zero.

For eventually nonzero $g$, write $f\sim g$ when
\[
f/g\to1.
\]

Equivalently $f=g+o(g)$, since
\[
(f-g)/g=f/g-1.
\]

Once the elementary sine limit is established, it says
$\sin x=x+o(x)$ and $\sin x\sim x$ as $x\to0$.

\paragraph{Preview: why differentiation uses little-o.}
Chapter~8's derivative condition is equivalent to
\[
f(a+h)=f(a)+f'(a)h+o(|h|)\qquad(h\to0).
\]
Indeed, for $r(h)=f(a+h)-f(a)-f'(a)h$, the absolute error quotient is

\[
|r(h)|/|h|=|(f(a+h)-f(a))/h-f'(a)|.
\]

Thus the nonlinear error is negligible compared with the first-order
scale, not merely a function tending to zero. The multivariable condition
$f(a+h)=f(a)+Df(a)h+o(\norm h)$ has exactly the same meaning.

\section{Basic Topology of the Real Line}
\label{sec:real-line-topology}

The local language used throughout analysis can be organized with a small
amount of topology.  We develop here only the real-line notions needed in
the remainder of Part~I; Chapter~\ref{ch:topology-linear-algebra} places them
in the general setting of metric and topological spaces.

\begin{definition}[Euclidean balls]
For $x\in\R$ and $r>0$, put
\[
B(x,r)=\set{y\in\R:|y-x|<r}=(x-r,x+r).
\]
\end{definition}

\begin{definition}[Intervals, open sets, and closed sets]
A set $I\subseteq\R$ is an \emph{interval} if $x,y\in I$ and $x<z<y$
imply $z\in I$.  A set $U\subseteq\R$ is \emph{open} if
\[
\forall x\in U\ \exists r>0:\quad B(x,r)\subseteq U.
\]
A set $F\subseteq\R$ is \emph{closed} if $\R\setminus F$ is open.
\end{definition}

An \emph{open neighborhood} of $x$ is an open subset of $\R$ containing
$x$. More generally, a \emph{neighborhood} of $x$ is a set containing an
open neighborhood of $x$.

The intervals $(a,b)$, $(-\infty,a)$, and $(a,\infty)$ are open, while
$[a,b]$, $(-\infty,a]$, and $[a,\infty)$ are closed.  Both $\varnothing$
and $\R$ are open and closed: the defining condition is vacuous for the
empty set, and every real point has a ball contained in $\R$.

\begin{proposition}[Elementary set operations]
Arbitrary unions and finite intersections of open subsets of $\R$ are
open.  Arbitrary intersections and finite unions of closed subsets of
$\R$ are closed.
\end{proposition}
\begin{proof}
If $x\in\bigcup_\alpha U_\alpha$, then $x\in U_\alpha$ for at least one
index, and a ball witnessing openness of that member lies in the union.
If $x\in U_1\cap\cdots\cap U_n$, choose radii $r_i>0$ with
$B(x,r_i)\subseteq U_i$ and set
\[
r=\min\{r_1,\ldots,r_n\}>0.
\]
Then $B(x,r)$ lies in every $U_i$.  The closed-set assertions follow by
taking complements and using De Morgan's laws.
\end{proof}

Finiteness matters: for example,
\[
\bigcap_{n=1}^\infty(-1/n,1/n)=\{0\},
\]
which is not open.  Taking complements also gives an infinite union of
closed sets that is not closed.

\begin{definition}[Relative openness and closedness]
Let $A\subseteq Y\subseteq\R$.  The set $A$ is \emph{relatively open in
$Y$} if $A=Y\cap U$ for some open $U\subseteq\R$.  It is
\emph{relatively closed in $Y$} if $A=Y\cap F$ for some closed
$F\subseteq\R$; equivalently, $Y\setminus A$ is relatively open in $Y$.
\end{definition}
For instance, if $a<c<b$, then
\[
[a,c)=[a,b]\cap(-\infty,c)
\]
is relatively open in $[a,b]$ although it is not open in $\R$.

Limit points and isolated points were defined in \cref{prop:limit-point-sequence}
and the paragraphs preceding it.  Those notions record whether every
punctured neighborhood meets a set; the sequential test there will remain
our principal tool.

\begin{theorem}[Sequential characterization of closed sets]
\label{thm:sequentially-closed}
A set $F\subseteq\R$ is closed if and only if every convergent sequence in
$F$ has its limit in $F$.
\end{theorem}
\begin{proof}
Suppose first that $F$ is closed and $x_n\in F$ converges to $x$. If
$x\notin F$, the open set $F^c$ contains a ball $B(x,r)$. Eventually
$x_n\in B(x,r)$, a contradiction. Conversely, suppose the sequential
condition holds and let $x\in F^c$. If no ball about $x$ lay in $F^c$,
then for every $n$ one could choose $x_n\in F$ with $|x_n-x|<1/n$.
This sequence converges to $x$, contradicting the hypothesis. Hence $F^c$
is open.
\end{proof}

\begin{definition}[Dense subsets]
For $E\subseteq Y\subseteq\R$, the set $E$ is \emph{dense in $Y$} if
every nonempty relatively open subset of $Y$ meets $E$. Equivalently, every
neighborhood of a point of the subspace $Y$ meets $E$.
\end{definition}
By \cref{cor:rationals-countable,thm:rational-density},
$\Q\cap[a,b]$ is countable and dense in $[a,b]$.

\begin{proposition}[Interval components of an open set]
\label{prop:real-open-interval-components}
Every open subset of $\R$ is a disjoint union of at most countably many
open intervals.
\end{proposition}
\begin{proof}
Let $G\subseteq\R$ be open.  For each $x\in G$, take the union of all open
intervals that contain $x$ and lie in $G$.  Overlapping intervals merge
into an interval, so this union is a maximal open interval contained in
$G$. Two such maximal intervals are either equal or disjoint, and together
they cover $G$. Fix an enumeration $\Q=\{q_1,q_2,\ldots\}$. Associate to
each component interval the first $q_n$ that it contains. Distinct disjoint
components receive distinct rationals, so this gives an injection of the
component intervals into $\N$.
\end{proof}

\section{Continuity}

Continuity compares nearby function values with the value at a domain point.
At a limit point it removes the separation between the limiting value and
the assigned value; isolated points require a separate explanation.  The
topology of \cref{sec:real-line-topology} gives a second language for this
stability.

\begin{definition}[Continuity]
\label{def:continuity}
Let $f:D\to\R$ and let $a\in D$. The function $f$ is \emph{continuous at $a$} if, for every $\eps>0$, there is $\delta>0$ such that every $x\in D$ satisfying $\abs{x-a}<\delta$ satisfies $\abs{f(x)-f(a)}<\eps$. It is \emph{continuous on $D$} if it is continuous at every point of $D$.
\end{definition}

\begin{theorem}[Open-set characterization of continuity]
\label{thm:real-open-set-continuity}
For $f:D\to\R$, the function $f$ is continuous on $D$ if and only if
$f^{-1}(U)$ is relatively open in $D$ for every open $U\subseteq\R$.
Equivalently, inverse images of closed sets are relatively closed in $D$.
\end{theorem}
This result is recorded to connect the $\eps$--$\delta$ definition with the
topological language used later. Its proof is not given here; Chapter~10
proves the corresponding statement for general topological spaces in
\cref{thm:topological-continuity}.

\paragraph{Limits and continuity have different domain requirements.}
A function limit requires $a$ to be a limit point but need not have
$a\in D$; its test $0<|x-a|<\delta$ excludes $x=a$. Continuity
requires $a\in D$ because $f(a)$ must exist, but need not require a
limit point; its test $|x-a|<\delta$ includes $a$.

For $D=\{0\}\cup[1,2]$, choose
\[
\delta=1/2.
\]
 Any $x\in D$ with
$|x|<\delta$ must equal $0$. The only comparison required for continuity
is therefore $|f(0)-f(0)|=0<\eps$, for every $\eps>0$, regardless
of the other values of $f$. Within that neighborhood the function is
defined only at the isolated point itself. This differs from a punctured
limit, where there would be no admissible points at all.

Consequently, if $a\in D$ is a limit point, $f$ is continuous at $a$
if and only if
\[
\lim_{x\to a}f(x)=f(a):
\]
 the excluded case $x=a$
adds only the comparison $0<\eps$. If $a\in D$ is isolated, every
function on $D$ is continuous there, while its punctured function limit
is not defined under the convention of these notes.

\begin{theorem}[Sequential criterion for continuity]
\label{thm:sequential-continuity}
Let $f:D\to\R$ and $a\in D$. Then $f$ is continuous at $a$ if and only if every sequence $(x_n)$ in $D$ satisfying $x_n\to a$ also satisfies $f(x_n)\to f(a)$.
\end{theorem}
\begin{proof}
The proof is the same as \cref{thm:sequential-function-limit}, except that terms equal to $a$ require no exclusion: for such a term, $\abs{f(x_n)-f(a)}=0$. Conversely, failure of continuity gives, for each $n$, an $x_n\in D$ with
\[
\abs{x_n-a}<1/n
\]
 and $\abs{f(x_n)-f(a)}\geq\eps_0$.
\end{proof}

\begin{theorem}[Algebra of continuous functions]
\label{thm:continuity-algebra}
Let $f,g:D\to\R$ be continuous at $a\in D$, and let $c\in\R$. Then $f+g$, $cf$, and $fg$ are continuous at $a$. If $g(a)\ne0$, then
\[
f/g
\]
 is continuous at $a$ on the set where it is defined.
\end{theorem}
\begin{proof}
For every sequence $x_n\to a$ in $D$, continuity gives $f(x_n)\to f(a)$ and $g(x_n)\to g(a)$. Apply the sequence limit laws of \cref{thm:limit-algebra} and then \cref{thm:sequential-continuity}.
\end{proof}

\begin{proposition}[Composition preserves continuity]
\label{prop:continuous-composition-real}
Let $D,E\subseteq\R$, let $f:D\to E$ and $g:E\to\R$, and let $a\in D$.
If $f$ is continuous at $a$ and $g$ is continuous at $f(a)$, then
$g\circ f$ is continuous at $a$.
\end{proposition}
\begin{proof}
For every sequence $x_n\to a$ in $D$, the sequential criterion
\cref{thm:sequential-continuity} gives $f(x_n)\to f(a)$ in $E$ and then
$g(f(x_n))\to g(f(a))$. The same criterion proves the conclusion.
\end{proof}

\begin{example}
Every polynomial is continuous on $\R$: constant and identity functions are continuous directly from \cref{def:continuity}, and repeated application of \cref{thm:continuity-algebra} proves the claim. Every rational function is continuous wherever its denominator is nonzero.
\end{example}

\begin{example}[Checking a junction and a removable value]
Let $f(x)=x+1$ for $x<0$ and $f(x)=x^2+c$ for $x\geq0$.
Away from zero use polynomial continuity. At zero the left limit is
$1$ and the right limit and assigned value are $c$, so continuity holds
exactly when $c=1$. For a removable example, assign a value $d$ at
$x=1$ to
\[
(x^2-1)/(x-1).
\]
 The punctured limit is $2$, so its
extension is continuous exactly when $d=2$. Altering the value cannot
repair a jump whose two side limits disagree.
\end{example}

\begin{definition}[Bounded functions]
\label{def:bounded-function}
A function $f:D\to\R$ is \emph{bounded} if
\[
\exists M\geq0\ \forall x\in D:\quad |f(x)|\leq M.
\]
Equivalently, its range $f(D)$ is a bounded subset of $\R$.
\end{definition}

\begin{definition}[Piecewise continuity]
\label{def:piecewise-continuous}
A function $f:[a,b]\to\R$ is \emph{piecewise continuous} if there is a
finite partition
\[
a=x_0<x_1<\cdots<x_n=b
\]
such that, on every $(x_{i-1},x_i)$, the restriction of $f$ extends
continuously to $[x_{i-1},x_i]$. Thus only finitely many breakpoints are
allowed, and finite one-sided limits exist there.
\end{definition}

\section{Uniform Continuity}

Ordinary continuity permits the required input tolerance to vary from point
to point.  Uniform continuity asks for one tolerance that works everywhere
on the domain, precisely the stronger control needed for approximation and
integration.

\begin{definition}[Uniform continuity]
\label{def:uniform-continuity}
Let $f:D\to\R$. The function $f$ is \emph{uniformly continuous on $D$} if, for every $\eps>0$, there is $\delta>0$ such that, for all $x,y\in D$,
\[
\abs{x-y}<\delta\Longrightarrow\abs{f(x)-f(y)}<\eps.
\]
\end{definition}

The difference from ordinary continuity is that $\delta$ may depend on $\eps$, but not on the point of $D$.

Schematically, the quantifiers are
\[
\text{continuity:}\quad
\forall a\in D\;\forall\eps>0\;\exists\delta=\delta(a,\eps)>0,
\]
whereas
\[
\text{uniform continuity:}\quad
\forall\eps>0\;\exists\delta=\delta(\eps)>0\;\forall x,y\in D.
\]
Moving the point quantifiers past $\exists\delta$ is the whole strengthening.

\begin{proposition}[Lipschitz continuity implies uniform continuity]
\label{prop:lipschitz-uniform}
Let $f:D\to\R$. If there is $K\geq0$ such that
\[
\abs{f(x)-f(y)}\leq K\abs{x-y}
\qquad\text{for all }x,y\in D,
\]
then $f$ is uniformly continuous on $D$.
\end{proposition}
\begin{proof}
If $K=0$, the function is constant. If $K>0$, choose
\[
\delta=\eps/K.
\]
 Then the displayed inequality gives the conclusion of \cref{def:uniform-continuity}.
\end{proof}

 A function satisfying the displayed inequality is \textbf{Lipschitz
continuous}, or simply \textbf{Lipschitz}, on (D), and (K) is a
\textbf{Lipschitz constant}. Thus the hierarchy is
\[
\text{Lipschitz continuity}\Longrightarrow\text{uniform continuity}
\Longrightarrow\text{continuity}.
\]
Neither converse holds in general.

\begin{proposition}[Sequence pairs detect failure of uniform continuity]
\label{prop:uniform-continuity-pairs}
A function $f:D\to\R$ is uniformly continuous if and only if every
pair of sequences $x_n,y_n\in D$ with $|x_n-y_n|\to0$ satisfies
$|f(x_n)-f(y_n)|\to0$.
\end{proposition}
\begin{proof}
For uniform continuity, given $\eps>0$ choose its common $\delta>0$.
Choose $N$ so that $|x_n-y_n|<\delta$ for $n\geq N$.
Then $|f(x_n)-f(y_n)|<\eps$ for every $n\geq N$.
Conversely, failure gives $\eps_0>0$ such that for each $n$ there are
$x_n,y_n\in D$ with
\[
|x_n-y_n|<1/n
\]
 but
$|f(x_n)-f(y_n)|\geq\eps_0$. These pairs contradict the stated test.
\end{proof}

\begin{example}[The three levels of control]
The function $x\mapsto x^2$ is not uniformly continuous on $\R$: with
$x_n=n$ and
\[
y_n=n+1/n,
\]
 the inputs approach one another while
\[
y_n^2-x_n^2=2+\frac1{n^2}.
\]
On $[-A,A]$, however,
\[
\abs{x^2-y^2}\leq2A\abs{x-y},
\]
so it is Lipschitz there.  The function $\sqrt{x}$ on $[0,1]$ is uniformly
continuous because
\[
\abs{\sqrt{x}-\sqrt{y}}\leq\sqrt{\abs{x-y}},
\]
as follows: if $u=\sqrt{x}\geq v=\sqrt{y}$, then

\[
(u-v)^2\leq(u-v)(u+v)=x-y;
\]
 interchange $x,y$ for the other order.
The uniform-continuity proof chooses $\delta=\eps^2$.
It is not Lipschitz: if a
constant $K$ worked, then $\sqrt{x}\leq Kx$ for every $x>0$, equivalently

\[
1/\sqrt{x}\leq K,
\]
 which fails near $0$.
\end{example}

\begin{example}[Continuity need not be uniform]
The function $f:(0,1)\to\R$ given by
\[
f(x)=\frac1x
\]
 is continuous by
\cref{thm:continuity-algebra}, but not uniformly continuous. Indeed, put
\[x_n=\frac1{n+1},\qquad y_n=\frac1{n+2}.\]
Then
$\abs{x_n-y_n}\to0$, whereas $\abs{f(x_n)-f(y_n)}=1$ for every $n$.
\end{example}

On $[a,\infty)$ with $a>0$, the same function
\[
1/x
\]
 is Lipschitz, since
\[
\abs{\frac1x-\frac1y}=\frac{\abs{x-y}}{xy}
\leq\frac1{a^2}\abs{x-y}.
\]
This contrast shows that uniform continuity depends on the domain as well
as on the formula.

\section{Intermediate Values}

Continuity has a qualitative consequence beyond limiting estimates: it
prevents a function from jumping over a value.  Completeness supplies the
point at which this intuitive crossing must occur.

\begin{theorem}[Intermediate value theorem]
\label{thm:ivt}
Let $f:[a,b]\to\R$ be continuous, with $a<b$. If $u$ lies between $f(a)$ and $f(b)$, then there is $c\in[a,b]$ such that $f(c)=u$.
\end{theorem}
\begin{proof}
It is enough to treat $f(a)\leq u\leq f(b)$. Let
\[
S=\set{x\in[a,b]:f(x)\leq u}.
\]
The set is nonempty and bounded above, so let $c=\sup S$. If $f(c)<u$, then $c\ne b$ because
\[
f(b)\geq u.
\]
Continuity gives a positive $\delta$ such that $f(x)<u$ whenever
\[
\abs{x-c}<\delta.
\]
Choose
\[
0<r<\min(\delta,b-c).
\]
Then $c+r\in S$, contradicting that $c$ is an upper bound of $S$. If $f(c)>u$, then $c\ne a$ because
\[
f(a)\leq u.
\]
Continuity gives a positive $\delta$ such that $f(x)>u$ whenever
\[
\abs{x-c}<\delta.
\]
By the defining property of the supremum, there is
\[
x\in S
\]
with
\[
c-\delta<x\leq c,
\]
which contradicts $f(x)\leq u$. Therefore $f(c)=u$.
\end{proof}

\begin{theorem}[Continuity of a monotone inverse]
\label{thm:monotone-inverse-continuity}
If $f$ is continuous and strictly monotone on an interval $I$, then
$f:I\to f(I)$ is bijective and its inverse is continuous.
\end{theorem}
\begin{proof}
Strict monotonicity gives injectivity, while surjectivity onto $f(I)$ is
automatic. The IVT shows that $f(I)$ is an interval. Suppose first that
$f$ is increasing, fix $a\in I$, and write $y_0=f(a)$.

If $a$ is an interior point of $I$, then for a given $\eps>0$ choose
\[
a-\eps<\ell<a<r<a+\eps,
\qquad \ell,r\in I.
\]
The number
\[
\delta=\min\{f(a)-f(\ell),\,f(r)-f(a)\}>0
\]
has the desired effect: if $y\in f(I)$ and $|y-y_0|<\delta$, then
\[
f(\ell)<y<f(r),
\]
so monotonicity traps $f^{-1}(y)$ between $\ell$ and $r$, and hence within
$\eps$ of $a$.

At a left endpoint $a$, choose only $r\in I$ with
$a<r<a+\eps$ and take $\delta=f(r)-f(a)$.  Then
$a\leq f^{-1}(y)<r$ whenever $y\in f(I)$ and $|y-y_0|<\delta$.
The right-endpoint argument is symmetric; a singleton interval is
immediate. Thus $f^{-1}$ is continuous at every point of $f(I)$. If $f$
is decreasing, apply the increasing case to $-f$. Since
$f^{-1}(y)=(-f)^{-1}(-y)$, continuity follows by composition with negation.
\end{proof}

Integer powers are already defined. We now reverse the question: given
$a$, can we solve $x^n=a$? Elementary algebra uses radical notation for a
chosen solution; analysis must justify its existence and uniqueness.

\begin{theorem}[Existence and uniqueness of roots]
\label{thm:real-roots}
For every $c\geq0$ and $n\in\N$, there is a unique $r\geq0$ such that
$r^n=c$.
\end{theorem}
\begin{proof}
The polynomial $p(t)=t^n-c$ is continuous.  On
$[0,\max\{1,c\}]$ its endpoint values have opposite weak signs, so the
Intermediate Value Theorem gives a zero.  If $0\leq r<s$, then
\[
s^n-r^n=(s-r)(s^{n-1}+s^{n-2}r+\cdots+r^{n-1})>0,
\]
so $t\mapsto t^n$ is strictly increasing on $[0,\infty)$ and the root is
unique.
\end{proof}

\begin{remark}[Roots of an equation and the principal radical]
Saying that $x$ is an $n$th root of $a$ means $x^n=a$.
The symbol $\sqrt[n]{a}$ instead chooses the \emph{principal real root}.
For example, if $a>0$, then
\[
x^2=a\quad\Longleftrightarrow\quad x=\pm\sqrt a,
\]
whereas $\sqrt a$ alone denotes the unique nonnegative solution.
For even $n$, the principal root is nonnegative and exists exactly when
$a\geq0$; for $a>0$ the equation also has its negative as a solution.
For odd $n$, every real $a$ has exactly one real root, negative when $a<0$,
given in that case by $-\sqrt[n]{-a}$. These assertions follow from the
root theorem and the parity of integer powers. Analysis has now justified
the familiar radicals, rather than assuming that writing a symbol produces
a solution.
\end{remark}

\begin{remark}[Continuity of roots]
Positive square roots were constructed directly in
\cref{prop:positive-square-roots}, in Chapter~\ref{ch:reals}.
The general positive $n$th roots established above using the IVT
also vary continuously. The map \(t\mapsto t^n\) is continuous and
strictly increasing on \([0,\infty)\). Its inverse
\(x\mapsto\sqrt[n]{x}\) is therefore continuous by
\cref{thm:monotone-inverse-continuity}. At the endpoint this also has the
explicit tolerance
\[
0\leq x<\eps^n\quad\Longrightarrow\quad
0\leq\sqrt[n]{x}<\eps.
\]
For odd roots on negative inputs, use the sign formula above and the same
conclusion on each side of zero.
\end{remark}

\begin{proposition}[Two-variable arithmetic--geometric mean inequality]
\label{prop:two-variable-amgm}
For \(a,b\geq0\),
\[
\boxed{\sqrt{ab}\leq\frac{a+b}{2}.}
\]
Equality holds exactly when \(a=b\).
\end{proposition}
\begin{proof}
The quadratic inequality in \cref{prop:quadratic-bernoulli} says
\(2uv\leq u^2+v^2\). Set \(u=\sqrt a\) and \(v=\sqrt b\), whose existence
has now been established. Then
\[
2\sqrt{ab}=2\sqrt a\sqrt b\leq a+b.
\]
Equality in the quadratic inequality holds exactly when
\(\sqrt a=\sqrt b\), equivalently \(a=b\).
\end{proof}

\begin{proposition}[Laws of real roots]
\label{prop:real-root-laws}
Let \(n\in\N\), let \(a,b\geq0\), and suppose \(b>0\) in the quotient
formula. Then
\[
\sqrt[n]{ab}=\sqrt[n]{a}\sqrt[n]{b},
\qquad
\sqrt[n]{\frac ab}=\frac{\sqrt[n]{a}}{\sqrt[n]{b}}.
\]
For odd \(n\), the product law also holds for arbitrary real \(a,b\) when
the unique real odd roots are used.
\end{proposition}
\begin{proof}
Both sides of the product identity are nonnegative, and raising the right
side to the \(n\)th power gives \(ab\). Uniqueness in
\cref{thm:real-roots} proves the equality. The quotient proof is identical:
the displayed quotient is nonnegative and its \(n\)th power is
\[
a/b.
\]

For odd \(n\), signs multiply in the same way, and uniqueness among all real
odd roots completes the argument.
\end{proof}

To extend integer exponent laws, the value
\[
a^{1/n}
\]
 for $a>0$ must
satisfy
\[
(a^{1/n})^n=a.
\]
 Requiring positive values therefore forces the
positive $n$th root. The same laws then suggest
\[
a^{1/n}=\sqrt[n]{a},\qquad
a^{m/n}=(a^{1/n})^m=\sqrt[n]{a^m}.
\]
The last equality follows by raising both positive expressions to the
$n$th power and using uniqueness. Thus the progression is from integer
powers to solving an equation, choosing its principal root, and extending
the exponent. We still have to check that different fractions representing
one rational number give the same value.

\begin{definition}[Positive rational powers]
\label{def:rational-powers}
Let \(a>0\) and \(r\in\Q\). Choose \(m\in\Z\) and \(n\in\N\) with
\[
r=\frac mn,
\]
and define
\[
\boxed{a^r=a^{m/n}:=\sqrt[n]{a^m}.}
\]
Integer powers with negative exponent have their usual reciprocal meaning.
The next theorem shows that the value is independent of the chosen fraction
representing \(r\).
\end{definition}

\begin{theorem}[Rational-exponent laws]
\label{lem:rational-power-laws}
For \(a,b>0\) and \(r,s\in\Q\), the preceding definition is well defined
and satisfies
\[
\begin{aligned}
a^{r+s}&=a^ra^s,&
a^{r-s}&=\frac{a^r}{a^s},\\
(a^r)^s&=a^{rs},&
(ab)^r&=a^rb^r.
\end{aligned}
\]
If \(0<a<b\) and \(r>0\), then \(a^r<b^r\). For fixed \(a>1\), the map
\(r\mapsto a^r\) is strictly increasing on \(\Q\); for \(0<a<1\), it is
strictly decreasing.
\end{theorem}
\begin{proof}
Suppose
\[
m/n=p/q,
\]
 with \(n,q>0\), and set
\(u=\sqrt[n]{a^m}\), \(v=\sqrt[q]{a^p}\). Both are positive, and
\[
u^{nq}=a^{mq}=a^{pn}=v^{nq}.
\]
Uniqueness of positive \(nq\)th roots gives \(u=v\). This is the promised
well-definedness: the notation depends on the rational number, not on how it
is written as a fraction.

For addition, put the exponents over one positive denominator,
say $r=m/n$ and $s=p/n$. Both sides below are positive, and
\[
(a^ra^s)^n=(a^r)^n(a^s)^n=a^ma^p=a^{m+p}
                 =(a^{r+s})^n.
\]
Uniqueness of the positive $n$th root therefore gives
$a^{r+s}=a^ra^s$.

For the iterated-power law, write $r=m/n$ and $s=p/q$, with $n,q>0$.
After raising both positive candidates to the common power $nq$, integer
exponent laws give
\[
\bigl((a^r)^s\bigr)^{nq}=a^{mp}
               =\bigl(a^{rs}\bigr)^{nq}.
\]
Uniqueness of the positive $nq$th root proves $(a^r)^s=a^{rs}$.
The quotient law follows from the addition law applied to $r-s+s=r$, and
the product law follows by raising both positive sides to a positive
denominator of $r$ and using \cref{prop:real-root-laws}.

If \(0<a<b\) and
\[
r=m/n>0,
\]
 then \(m,n\in\N\) may be chosen positive.
Thus \(a^m<b^m\). Were \(a^r\geq b^r\), raising to the \(n\)th power would
give \(a^m\geq b^m\), contradicting \(a^m<b^m\). Finally,

\[
a^s/a^r=a^{s-r}>1
\]
 when \(a>1\) and \(s>r\); taking reciprocals gives
the statement for \(0<a<1\).
\end{proof}

\begin{remark}[Why the clean theory uses positive bases]
The equality
\[
\frac13=\frac26
\]
 exposes a representation problem for
negative bases. Interpreting each fraction through a principal radical gives
\[
\sqrt[3]{-8}=-2,\qquad \sqrt[6]{(-8)^2}=2.
\]
Thus the same rational exponent would give incompatible answers under a
naive representation-by-representation rule.
Some negative-base expressions are meaningful: for example,

\[
(-8)^{1/3}=-2,
\]
 and more generally odd roots of negative numbers exist.
Even roots of negative numbers are not real, however, and a reduced rational
exponent may impose a parity restriction on its denominator. To avoid hiding
these domain issues, the general rational-power theory is formulated for
positive bases. Chapter~\ref{ch:differentiation} will extend it to arbitrary
real exponents, still for positive bases, by means of \(\exp\) and \(\log\).
\end{remark}

\section{Compactness on the Real Line}
\label{sec:real-line-compactness}

Compactness is the mechanism that turns local information into finitely
many choices and one common scale. We record both formulations used in
Part~I.

\begin{definition}[Covers and subcovers]
Let $K\subseteq\R$. A family $\{U_\alpha\}_{\alpha\in A}$ \emph{covers}
$K$ if $K\subseteq\bigcup_{\alpha\in A}U_\alpha$. It is an \emph{open
cover} if every $U_\alpha$ is open in $\R$. If $B\subseteq A$ and the
members indexed by $B$ still cover $K$, they form a \emph{subcover}; it is
a \emph{finite subcover} when $B$ is finite.
\end{definition}
A subcover selects members of the original family; it does not replace
them by newly constructed sets. The cover condition is containment, not
equality: the single interval $(-1,2)$ covers $[0,1]$.

\begin{definition}[Compactness and sequential compactness]
\label{def:real-compactness}
A set $K\subseteq\R$ is \emph{compact} if every open cover of $K$ has a
finite subcover. It is \emph{sequentially compact} if every sequence in
$K$ has a subsequence converging to an element of $K$.
\end{definition}

\begin{theorem}[Compactness and sequential compactness on the line]
\label{thm:real-compact-sequential}
A subset of $\R$ is compact if and only if it is sequentially compact.
\end{theorem}
The general equivalence is not proved here. Part~I develops the two
formulations directly in the concrete settings it needs; Chapter~10 proves
the corresponding metric-space result in
\cref{thm:metric-compact-sequential}.

\begin{theorem}[Sequential compactness of closed bounded intervals]
\label{thm:compact-interval}
If $a\leq b$, then $[a,b]$ is sequentially compact.
\end{theorem}
\begin{proof}
Let $(x_n)$ be a sequence in $[a,b]$. It is bounded, so \cref{thm:bolzano-weierstrass} gives a subsequence converging to some $x\in\R$. Since $a\leq x_{n_k}\leq b$ for every $k$, \cref{thm:limit-order} gives $a\leq x\leq b$. Thus $x\in[a,b]$.
\end{proof}

For an interval $J$, write $\ell(J)$ for its ordinary length.

\begin{lemma}[Common scale for an interval cover]
\label{lem:lebesgue-number-interval}
For every open cover $\mathcal U$ of $[a,b]$, there is $\lambda>0$ such
that every subinterval of $[a,b]$ having length less than $\lambda$ is
contained in one member of $\mathcal U$.
\end{lemma}
\begin{proof}
Suppose no such $\lambda$ exists. For each $n\in\N$, choose an interval
$J_n\subseteq[a,b]$ of length less than $1/n$ that lies in no member of
$\mathcal U$, and choose $x_n\in J_n$. By
\cref{thm:compact-interval}, some subsequence $x_{n_k}$ converges to
$x\in[a,b]$. Choose $U\in\mathcal U$ containing $x$ and $r>0$ with
$B(x,r)\subseteq U$. For all sufficiently large $k$,
\[
|x_{n_k}-x|<r/2,
\qquad
\ell(J_{n_k})<1/n_k<r/2.
\]
Every $y\in J_{n_k}$ then satisfies
$|y-x|\leq|y-x_{n_k}|+|x_{n_k}-x|<r$, so $J_{n_k}\subseteq U$, a
contradiction.
\end{proof}

\begin{theorem}[Compactness of closed intervals]
\label{thm:compact-closed-interval}
Every closed bounded interval $[a,b]$ is compact.
\end{theorem}
\begin{proof}
Given an open cover, take $\lambda$ from
\cref{lem:lebesgue-number-interval}. Choose a finite partition of
$[a,b]$ whose cells have length less than $\lambda$. Each cell lies in one
member of the cover, and selecting one such member for each of the finitely
many cells gives a finite subcover.
\end{proof}

\begin{corollary}[Closed subsets of compact intervals]
\label{cor:closed-subset-finite-subcover}
If $K\subseteq[a,b]$ is relatively closed in $[a,b]$, then $K$ is compact.
\end{corollary}
\begin{proof}
Let $\mathcal U$ be an open cover of $K$. Relative closedness makes
$[a,b]\setminus K$ relatively open, so it equals $[a,b]\cap W$ for some
open $W\subseteq\R$. Then $\mathcal U\cup\{W\}$ covers $[a,b]$ and has a
finite subcover by \cref{thm:compact-closed-interval}. Discarding $W$, if
it was selected, leaves finitely many original members that cover $K$.
\end{proof}

Boundedness of the domain supplies a convergent subsequence of trial
points, closedness keeps its limit in the domain, and continuity transfers
their function values to that limit. All three steps matter: $x$ on
$(0,1)$ approaches its supremum without attaining it, while
\[
1/x
\]
 on
$(0,1]$ is continuous but unbounded.

\begin{theorem}[Extreme value theorem]
\label{thm:extreme-value}
Let $f:[a,b]\to\R$ be continuous, where $a\leq b$. Then there are points $x_{\min},x_{\max}\in[a,b]$ such that
\[
f(x_{\min})\leq f(x)\leq f(x_{\max})
\qquad\text{for every }x\in[a,b].
\]
\end{theorem}
\begin{proof}
Let $s=\sup f([a,b])$. To justify that this supremum exists, suppose first that $f([a,b])$ were unbounded above. For each $n$, choose $x_n\in[a,b]$ with $f(x_n)>n$. By \cref{thm:compact-interval}, a subsequence $x_{n_k}$ converges to $x\in[a,b]$. Continuity gives $f(x_{n_k})\to f(x)$, but $f(x_{n_k})\geq n_k\to\infty$, an impossibility. Hence the image is bounded above, and it is nonempty. For each $n$, choose $x_n\in[a,b]$ with
\[
f(x_n)>s-1/n.
\]
 A convergent subsequence has limit $x_{\max}\in[a,b]$. Continuity and the squeeze theorem give $f(x_{\max})=s$. Applying the same argument to $-f$ gives $x_{\min}$.
\end{proof}

\begin{theorem}[Heine--Cantor theorem]
\label{thm:heine-cantor}
Every continuous function $f:[a,b]\to\R$, where $a\leq b$, is uniformly continuous.
\end{theorem}
\begin{proof}
Suppose not. Then some $\eps_0>0$ has the property that, for every $n$, there are $x_n,y_n\in[a,b]$ such that
\[
\abs{x_n-y_n}<\frac1n
\qquad\text{and}\qquad
\abs{f(x_n)-f(y_n)}\geq\eps_0.
\]
By \cref{thm:compact-interval}, a subsequence $x_{n_k}$ converges to $x\in[a,b]$. Retain those same indices in the paired sequence. Then
\[
|y_{n_k}-x|\leq|y_{n_k}-x_{n_k}|+|x_{n_k}-x|
<\frac1{n_k}+|x_{n_k}-x|\longrightarrow0.
\]
Thus $y_{n_k}\to x$ as well. Choosing a different subsequence of the $y_n$
would lose the paired estimate, which concerns $x_n$ and $y_n$ at one index. Continuity gives both $f(x_{n_k})\to f(x)$ and $f(y_{n_k})\to f(x)$, so their difference tends to $0$, contradicting the second inequality.
\end{proof}



\subsection*{Roots, tolerances, and existence in use}
The continuous roots constructed near \cref{thm:real-roots} now join
our earlier continuity tools in concrete existence and tolerance problems.

\begin{example}[A threshold and a sensitivity guarantee]
A tank starts empty and has volume $V(t)=t+t^2$ litres after $t$
minutes, for $0\leq t\leq2$; this formula is the stated model.
The varying input is time and the target is $V(t)=3$.
The polynomial $V(t)-3$ is continuous, negative at $1$ and positive
at $2$. The IVT guarantees a filling time in $(1,2)$ minutes;
it does not by itself compute that time. Since

\[
V(y)-V(x)=(y-x)(1+x+y)>0
\]
 for $0\leq x<y\leq2$, it is unique.
Moreover $|V(t)-V(s)|\leq5|t-s|$ on this operating interval.
An input error below
\[
\eps/5
\]
 minutes therefore produces an output
error below $\eps$ litres, uniformly across the interval. If only
continuity were known, Heine--Cantor would still supply one tolerance,
but without this numerical value. Chapter~\ref{ch:differentiation}
will obtain Lipschitz constants from derivative bounds.
\end{example}

\section*{Exercises}
\addcontentsline{toc}{section}{Exercises}

\begin{hexercise}{fun-direct}
Prove directly that $5x+2\to7$ at $x=1$, $x^2\to9$ at $x=3$, and
\[
1/(x+1)\to1/2
\]
 at $x=1$. Display all auxiliary bounds and positive $\delta$ choices.
\end{hexercise}

\begin{hexercise}{fun-infinity}
Find the limits at both $+\infty$ and $-\infty$ of
\[
(2x^2+1)/(x^3-1),
\]

\[
(3x^2-x)/(x^2+2),,\qquad (-x^3+1)/(2x+3).
\]
 For one infinite limit give the complete target-and-cutoff argument.
\end{hexercise}

\begin{hexercise}{fun-removable}
Find all $c$ making the function
\[
(x^2-4)/(x-2)
\]
 for $x\ne2$, with value $c$ at $2$, continuous. Compare with
\[
x/|x|
\]
 at zero: can any assigned value remove its discontinuity?
\end{hexercise}

\begin{hexercise}{fun-pairs}
Prove $x^2$ is uniformly continuous on $[0,4]$ but not on $[0,\infty)$. Prove
\[
1/x
\]
 is Lipschitz on $[2,\infty)$ and determine whether it is uniformly continuous on $(0,2]$.
\end{hexercise}

\begin{hexercise}{fun-ivt-model}
A continuous temperature reading rises from $18$ degrees at time zero to $24$ degrees after one hour. Must it equal $21$ degrees at some time? Formulate the function, domain, and theorem. Must there be exactly one such time?
\end{hexercise}

\begin{hexercise}{fun-sensitivity}
A device maps an input length $x\in[1,3]$ centimetres to an output area $x^2$ square centimetres. Find one input tolerance ensuring output error below $0.06$ square centimetre for every pair of admissible inputs. Explain why the tolerance is uniform.
\end{hexercise}

\begin{exercise}
Using only integer powers and root notation, prove continuity of $x\mapsto\sqrt[3]{x}$ on $[0,\infty)$. Show that it is not Lipschitz near zero.
\end{exercise}

\begin{exercise}
Give continuous functions showing that the compact-interval hypotheses in the EVT and Heine--Cantor theorem cannot simply be dropped. Specify which conclusion fails in each example.
\end{exercise}

\begin{exercise}
For
\[
D=\{0\}\cup\{1/n:n\in\N\}\cup[2,3),
\]
 determine which of
$0,1,2,3$ belong to $D$, are limit points of $D$, or are isolated in $D$.
At which of these points is every function on $D$ automatically continuous?
\end{exercise}
\begin{exercise}
As $h\to0$, classify $\sqrt{|h|},|h|,|h|^2$ according to whether they
are $o(1)$, $O(|h|)$, or $o(|h|)$. Translate each answer into a bound
or quotient limit. Repeat the relative-size comparison of $x$ with
$x^2$ as $x\to\infty$ and explain why the remainder does not vanish.
\end{exercise}
\begin{exercise}
Let $r=o(g)$ and $s=O(g)$ in one fixed regime. Prove $r+s=O(g)$,
and give an example showing that $r+s=o(g)$ need not hold.
\end{exercise}

\begin{exercise}
Use the sequential criterion to prove that
\[
f(x)=x^2
\]
is continuous at every real number.
\end{exercise}

\begin{hexercise}{jump-continuity}
Let
\[
f(x)=
\begin{cases}
0,&x\leq0,\\
1,&x>0.
\end{cases}
\]
Determine all points at which $f$ is continuous, and explain your answer using sequences.
\end{hexercise}

\begin{exercise}
Let $f:[a,b]\to\R$ be continuous and suppose
\[
f(a)f(b)<0.
\]
Use the Intermediate Value Theorem to prove that $f$ has a zero in $(a,b)$.
\end{exercise}
